% REPORT.tex -- a fuller report on the Isabelle/HOL development.
% Compile with pdflatex (three passes for cross-references).

\documentclass[11pt]{article}

\usepackage[T1]{fontenc}
\usepackage[margin=1.1in]{geometry}
\usepackage{amsmath,amssymb,amsthm}
\usepackage{array}
\usepackage{booktabs}
\usepackage{longtable}
\usepackage[hidelinks]{hyperref}

\theoremstyle{plain}
\newtheorem{theorem}{Theorem}[section]
\newtheorem{proposition}[theorem]{Proposition}
\newtheorem{corollary}[theorem]{Corollary}
\newtheorem*{theorem*}{Theorem}

% Long identifiers must be breakable, or they overflow the measure.
\newcommand{\ubrk}{\_\discretionary{}{}{}}
\newcommand{\brk}{\discretionary{}{}{}}
\setlength{\emergencystretch}{3em}
\newcommand{\isa}[1]{\textup{\texttt{\small #1}}}
\newcommand{\thy}[1]{\textup{\texttt{\small #1}}}

\title{Dependency and Scope, formalised\\
  \large A report on the Isabelle/HOL development}
\author{}
\date{11 September 2026}

\begin{document}
\maketitle

\begin{abstract}
\noindent
This report describes an Isabelle/HOL formalisation of Halvorson's
\emph{Dependency and Scope: On the Translation Between Lemmon and Fitch
Proofs}. The development is 32{,}372 lines across 49 theories, with
1{,}234 lemmas, theorems and corollaries, and no admitted proofs. It is built
in two layers: a compact reconstruction of the calculi, in which the
mathematics is developed, and an exact layer reproducing the datatypes and
twenty-one rule checks of the \emph{How Logic Works} proof checker, against
which the paper's algorithms were implemented. The principal positive result is
a construction theorem at the exact types: every nonempty correct Lemmon proof
is translated by the checked translation into a Fitch proof that passes every
check, carries the same conclusion, and needs no additional premise. The
principal negative results are Theorem~10 and Conjectures~27 and~28, each
refuted at both layers. Along the way the formalisation corrected four claims
in the paper, identified fourteen defects in the supplied Haskell, two of which
affect soundness, and isolated two design decisions in the original checker
on which the main theorem turns out to depend.
\end{abstract}

\tableofcontents

\section{Introduction}

Lemmon and Fitch notations present natural deduction differently. A Lemmon
proof is a numbered list; each line carries a formula, a justification citing
earlier lines, and an explicit set of the assumptions it depends on. A Fitch
proof is a nested structure; assumptions open boxes, and the discharge of an
assumption closes one. The paper asks how the two are related: whether one can
always be translated into the other, whether the translation can be made
line-by-line, and what the dependency column is for.

Formalising a paper of this kind involves a choice that is worth stating
plainly at the outset. One can formalise the \emph{mathematics} --- reconstruct
the two calculi in whatever form is most convenient and prove the theorems
about them --- or one can formalise the \emph{implementation} --- reproduce the
concrete rule set of the checker the paper's algorithms were written against,
and prove the theorems about that. These are different obligations, and a
theorem about the first is not a theorem about the second. The rule rosters
differ, the correctness predicates differ, and, as \S\ref{sec:bridge} records,
the two do not embed in one another.

This development does both, in two layers:

\begin{description}
\item[The compact layer] reconstructs the calculi as \isa{fm}, \isa{just},
  \isa{pline}, \isa{fitch\ubrk{}item}. It is where the mathematics is developed:
  the erasure, the unfolding construction, the impossibility results.
\item[The exact layer] reproduces the datatypes and the twenty-one rule checks
  of the \emph{How Logic Works} checker under \isa{HL\ubrk{}}/\isa{hl} names,
  including its representation choices down to the empty-string marker for a
  vacuous quantifier. This is the authoritative specification for the
  executable system, and the layer in which the final results are stated.
\end{description}

\noindent
The report is organised by subject rather than by layer.
\S\ref{sec:calculi} fixes the objects and \S\ref{sec:correct} the several
notions of correctness, which are not interchangeable.
\S\ref{sec:semantics} covers soundness, \S\ref{sec:erasure} the erasure and the
dependency column, \S\ref{sec:direct} the line-by-line translation,
\S\ref{sec:construction} the unfolding construction in the compact layer and
\S\ref{sec:exact} at the exact types, which is the bulk of the recent work.
\S\ref{sec:impossibility} gives the negative results. \S\ref{sec:found} collects
what the formalisation found: corrections to the paper, defects in the
supplied code, and two observations about the checker's design.
\S\ref{sec:validation} describes how the development is checked and
\S\ref{sec:limits} what it does not establish.

\section{The two calculi}\label{sec:calculi}

\subsection{Lemmon proofs}

A Lemmon line is a quadruple: a dependency set, a line number, a formula, and a
justification. Writing the dependency set on the left, as the paper does, a
proof of $Q \to (P \to (P \wedge Q))$ looks like this.

\begin{center}
\begin{tabular}{@{}llll@{}}
\toprule
deps & line & formula & justification \\
\midrule
1    & (1) & $P$                          & A \\
2    & (2) & $Q$                          & A \\
1,2  & (3) & $P \wedge Q$                 & $\wedge$I 1,2 \\
2    & (4) & $P \to (P \wedge Q)$         & CP 1,3 \\
     & (5) & $Q \to (P \to (P \wedge Q))$ & CP 2,4 \\
\bottomrule
\end{tabular}
\end{center}

\noindent
Line 4 discharges assumption 1 --- so 1 leaves the dependency set --- and line 5
discharges assumption 2. The justification of a discharging rule names two line
numbers: the assumption being discharged, and the line at which the discharge
occurs.

In Isabelle the compact form is \isa{pline} with \isa{just} in
\thy{LF\ubrk{}Lemmon.thy}; the exact form is \isa{hl\ubrk{}line} with
\isa{hl\ubrk{}justification} in \thy{LF\ubrk{}HLW.thy}. The exact justification type has
twenty-one constructors, matching the checker: modus ponens and modus tollens,
double negation, conditional proof, reductio, the conjunction, disjunction and
biconditional rules, the four quantifier rules, the two equality rules,
excluded middle, quantifier negation, and an arbitrary propositional-consequence
rule \isa{HL\ubrk{}PropTaut}.

\subsection{Fitch proofs}

A Fitch proof is a list of items, each either a line or a subproof; a subproof
carries an assumption and a body which is again such a list. Compact
\isa{fitch\ubrk{}item} and \isa{subproof} live in \thy{LF\ubrk{}Fitch.thy}, exact
\isa{hl\ubrk{}fitch\ubrk{}item} and \isa{hl\ubrk{}subproof} in \thy{LF\ubrk{}HLW\ubrk{}Fitch.thy}.

Two structural facts about the exact form matter later. Flattening a Fitch
proof, \isa{hlFlattenFitch}, produces a list of triples --- number, formula,
rule --- with the boxes' contents spliced in place, and a subproof contributes
its assumption line first. And the well-formedness checker records, for each
line, which completed boxes are citable there: a discharge may name a pair
$(a,c)$ only when a box with head $a$ and last line $c$ has \emph{closed at
that line's own level}. That is \isa{hlFitchNestingFrom}, and
\S\ref{sec:impossibility} shows it is the engine of all three impossibility
results.

\subsection{The erasure}

The translation from Fitch to Lemmon is the easy direction, and the paper's
Definition~3 gives it: flatten the proof and compute each line's dependency set
from the rule and the sets of the lines it cites. This is \isa{delta} in the
compact layer and $\delta_H$ (\isa{hlFitchToLemmon}) in the exact one. It is a
total function on all Fitch objects, well-formed or not.

\section{What ``correct'' means}\label{sec:correct}

The development distinguishes three predicates on exact Lemmon proofs, and the
distinction is load-bearing rather than bureaucratic.

\begin{description}
\item[\isa{hlCorrect}] is the faithful reproduction of the supplied checker:
  every line passes \isa{hlStructureOK} and \isa{hlRuleOK}. It is deliberately
  \emph{not} repaired, because its purpose is to be what the Haskell does. It
  permits physically unsorted line lists and does not constrain the dependency
  set of an excluded-middle line.
\item[\isa{hlVerifiedCorrect}] adds dependency closure and a canonical
  ordering: positive, strictly increasing line numbers.
\item[\isa{hlPaperCorrect}] adds the paper's exact dependency arithmetic, in
  particular that an excluded-middle line depends on nothing
  (\isa{hlPaperCorrect\ubrk{}LEM\ubrk{}empty}).
\end{description}

\noindent
These are strictly ordered, with \isa{hlPaperCorrect\ubrk{}verified} the easy
inclusion. The reason to keep three is that the paper's claims about the
dependency column are true of the third and false of the second. The sharpest
statement of that is in \thy{LF\ubrk{}HLW\ubrk{}Canonical.thy}:

\begin{theorem*}[\isa{hlVerifiedCorrect\ubrk{}has\ubrk{}no\ubrk{}dependency\ubrk{}reconstructor}]
There is no function that recovers the dependency sets of every
\isa{hlVerifiedCorrect} proof from its stripped lines.
\end{theorem*}

\noindent
It is proved by exhibiting two accepted proofs with identical stripped data and
different dependency sets, so no function could send the one to both
(\isa{hlVerifiedCorrect\ubrk{}does\ubrk{}not\ubrk{}determine\ubrk{}dependencies}). This is what
forces the third predicate to exist.

On the Fitch side there is a similar layering: \isa{hlFitchWellFormed} is the
faithful reproduction of the checker's structural test, \isa{hlFitchVerified}
adds genuine nesting, premise closure, a top-level conclusion and correctness of
the Lemmon erasure, and \isa{hlFitchCorrect} adds the scope-based rule check.
\S\ref{sec:found} explains why the faithful version had to be left alone and the
repairs put in the layers above it.

\section{Semantics and soundness}\label{sec:semantics}

Both layers have a semantics and a soundness theorem, so that ``correct'' is
answerable to something other than itself.

For the compact calculus, \thy{LF\ubrk{}Semantics.thy} proves
\isa{lemmon\ubrk{}semantic\ubrk{}soundness} and \isa{fitch\ubrk{}semantic\ubrk{}soundness}: an
accepted proof's conclusion is a semantic consequence of the assumptions it
displays. \isa{derivOK\ubrk{}sound} does the same for derivation trees.

For the exact calculus, \thy{LF\ubrk{}HLW\ubrk{}Semantics.thy} --- 2{,}814 lines, the
largest theory in the development --- proves \isa{hlRuleOK\ubrk{}sound}: whenever a
line passes the checker's rule test, its formula follows from the formulas of
the lines it cites. This is done rule by rule, and the quantifier rules are
where the work is. \isa{hlForallIntroduction\ubrk{}sound} and
\isa{hlExistentialElimination\ubrk{}sound} require the eigenconstant conditions to be
read correctly, and both rest on soundness results for the checker's witness
inference (\isa{hlWitnessLists\ubrk{}forall\ubrk{}sound},
\isa{hlWitnessLists\ubrk{}exists\ubrk{}sound}) which say that an inferred witness list
really does describe an instantiation. \isa{hlPropositionalConsequence\ubrk{}sound}
covers the tautology rule by way of valuations, and
\isa{hlQuantifierNegationEquivalent\ubrk{}sound} the quantifier-negation rule.
\isa{hlVerifiedCorrect\ubrk{}line\ubrk{}sound} assembles these into the statement for
whole proofs.

\thy{LF\ubrk{}HLW\ubrk{}Fitch\ubrk{}Semantics.thy} proves \isa{hlFitchVerified\ubrk{}sound} and
\isa{hlFitchCorrect\ubrk{}sound} from the displayed outer premises. These use the
premise-closure check explicitly; they do not show it redundant.

\section{The erasure and the dependency column}\label{sec:erasure}

\paragraph{Theorem 4.} The paper states that $\delta$ is total and that a Fitch
proof is correct if and only if its erasure is. Totality holds by construction.
The forward implication is \isa{theorem\ubrk{}4\ubrk{}forward} in \thy{LF\ubrk{}Delta.thy},
proved from a comparison of Fitch scope with reconstructed dependencies. The
converse is \emph{false}, and the witness is the paper's own Theorem~22
example: erasure discards excess scope, and with it the eigenconstant violation
that made the Fitch proof incorrect. This is
\isa{theorem\ubrk{}4\ubrk{}backward\ubrk{}fails} in the compact layer and
\isa{hl\ubrk{}theorem\ubrk{}4\ubrk{}backward\ubrk{}fails} in the exact one.

\paragraph{Propositions 5 and 7, Corollary 6.} \isa{proposition\ubrk{}5} establishes
that reconstructed dependencies lie within scope, and may lie properly within
it; \isa{corollary\ubrk{}6} exhibits two distinct correct Fitch proofs with the same
erasure, so $\delta$ is not injective. \isa{proposition\ubrk{}7} is the statement
that the dependency column is redundant: from a correct proof with its
dependency sets deleted, they can be recomputed. It holds for compact
\isa{lemmonCorrect} and, as \isa{hl\ubrk{}proposition\ubrk{}7}, for exact
\isa{hlPaperCorrect}, with \isa{hlPaperCorrect\ubrk{}dependencies\ubrk{}unique} giving
uniqueness. As noted in \S\ref{sec:correct}, it fails for the weaker predicate,
and that failure is proved rather than assumed.

\section{The line-by-line translation}\label{sec:direct}

The paper's Definition~9 asks when a Fitch proof is a \emph{positional image}
of a Lemmon proof: same lines, in the same order, with corresponding rules.
\isa{positionalImage} in \thy{LF\ubrk{}Conjecture27.thy} formalises it, with
\isa{lineMatches} pinning each Fitch line's rule to the corresponding source
justification. Assumptions may correspond either to Fitch premises or to box
assumptions --- without that latitude the notion would be vacuously
unsatisfiable.

The direct algorithm \isa{lemmonToFitchDirect} walks the source once and opens
a box at each assumption; it is analysed in \thy{LF\ubrk{}Positional.thy}
(2{,}014 lines, 102 results). The main results establish what its output looks like when
it succeeds: sorted numbering, premises first, the right citations, a top-level
final line, and, by
\isa{lemmonToFitchDirect\ubrk{}fitchWF\ubrk{}iff\ubrk{}citations}, that its
well-formedness reduces exactly to a citation condition.
\isa{lemmonToFitchDirect\ubrk{}positional} in \thy{LF\ubrk{}Direct.thy} shows that when it
returns a proof, that proof is a positional image of the source.

The algorithm is guarded by seven checks. One of them, \isa{boxHeadError},
exists because of a defect: without it the algorithm can emit two boxes sharing
a last line, which is exactly the configuration \S\ref{sec:impossibility} shows
no well-formed Fitch proof can contain. \isa{checks\ubrk{}reject\ubrk{}only\ubrk{}imageless\ubrk{}sources}
records the partial converse that was established --- all six witnesses the
checks reject are genuinely imageless --- while being explicit that six
witnesses are not all sources. Whether the seven checks are \emph{necessary} as
well as sufficient is recorded as an open question rather than claimed.

\section{The unfolding construction, compact}\label{sec:construction}

When no positional image exists, the paper's Definition~21 gives a different
route: unfold the Lemmon proof into a derivation tree rooted at its conclusion,
then emit a Fitch proof from the tree. Four obligations follow, which the paper
labels L1--L4.

\begin{description}
\item[L1] Unfolding terminates on every correct nonempty source.
  \isa{unfolding\ubrk{}terminates} and \isa{L1} in \thy{LF\ubrk{}Unfold.thy}. The
  recursion is well-founded because citations point to smaller numbers, which
  is \isa{proposition\ubrk{}18}.
\item[L2, L3] Citations precede their uses and are in scope; premises are
  emitted once, at the top. \isa{L2\ubrk{}L3} in \thy{LF\ubrk{}WellFormed.thy}
  (1{,}906 lines).
\item[L4] Each rule transfers: the emitted line satisfies the corresponding
  Fitch rule, including the subproof endpoints the discharging rules need and
  the eigenconstant conditions. \isa{L4\ubrk{}derivation} in \thy{LF\ubrk{}Check.thy}
  (2{,}465 lines).
\end{description}

\noindent
Two supporting results are the paper's Lemma~23 and Corollary~24: renaming a
constant to a globally fresh one preserves a derivation (\isa{renaming}), and
leaves the open assumptions unchanged when the old constant does not occur in
them (\isa{renaming\ubrk{}open}). These are what make the eigenconstant repair
legitimate.

The construction theorem for the compact calculus is \isa{conjecture\ubrk{}26} in
\thy{LF\ubrk{}Conjecture.thy}, with \isa{translation\ubrk{}total} and the sequent
corollary \isa{conjecture\ubrk{}26\ubrk{}sequent}. It is a \emph{corrected} replacement
for the paper's Conjecture~26 in one respect: the original repairs only
universal introduction, and the repair must also be applied to existential
elimination, whose witness carries the same kind of freshness condition. This is
discussed in \S\ref{sec:found}.

\section{The construction at the exact types}\label{sec:exact}

Everything in \S\ref{sec:construction} is about the compact calculus. The
corresponding statements for the exact checker do not follow from it --- the
datatypes, the rule roster and the emitter are all different --- and
establishing them is the bulk of the recent work. It proceeds in five steps.

\subsection{Unfolding}
\isa{hlUnfoldDerivation\ubrk{}succeeds} supplies sufficient fuel by citation rank and
\isa{hlToDerivation\ubrk{}total} gives successful unfolding for every nonempty
verified source; \isa{hlToDerivation\ubrk{}correct} then proves the resulting tree
correct, with the same conclusion and no additional premises, for all
twenty-one rules (\isa{hlVerifiedCorrect\ubrk{}toDerivation},
\isa{hlPaperCorrect\ubrk{}toDerivation}).

\subsection{Renaming}
\isa{hlDerivationOK\ubrk{}rename} proves that renaming to a fresh constant preserves
every exact derivation rule, quantifier witness selection, equality and
propositional consequence included; \isa{hlRenameDerivation\ubrk{}open\ubrk{}unchanged}
is the Corollary~24 analogue. The hypotheses exclude the checker's empty-name
sentinel, which is not decoration: renaming that marker would corrupt the
witness-inference logic.

\subsection{Layout}
\thy{LF\ubrk{}HLW\ubrk{}Layout.thy} proves the emitter allocates a consecutive interval of
line numbers (\isa{hlAllocated}), \thy{LF\ubrk{}HLW\ubrk{}Fuel.thy} that the fuel supplied
is always sufficient, and \thy{LF\ubrk{}HLW\ubrk{}Conclusion.thy} that the emitted proof
concludes with the tree's formula.

\subsection{Structural well-formedness}
\thy{LF\ubrk{}HLW\ubrk{}Nesting.thy} proves \isa{hlEmitDerivationFuel\ubrk{}nesting} by
well-founded induction on derivation size across all twenty-one rules: the
emitted fragment satisfies the genuine nesting predicate, so every discharge
names a box that really closed at the right level. Ordinary citation visibility
then needs no second induction: \isa{hlFitchNestingFrom\ubrk{}scope} derives it,
because the scope check is the weaker of the two --- after a box it makes the
box's assumption and last line visible, which nesting does not, and it inspects
no discharge pairs at all.

\subsection{Rule transfer through the erasure}
This is the long pole. What must be shown is that erasing the emitted proof
gives a \emph{correct Lemmon proof}, which means each of the three checks in
\isa{hlVerifiedCorrect}:

\begin{itemize}
\item \isa{hlFitchToLemmon\ubrk{}canonicalOrder} --- the erased numbers are the
  flattened Fitch numbers, allocated strictly increasingly;
\item \isa{hlFitchToLemmon\ubrk{}dependencyClosed} --- which holds of \emph{every}
  Fitch proof, with no hypothesis, because \isa{hlFitchDependenciesOf} computes
  exactly the union each rule demands;
\item \isa{hlFitchToLemmon\ubrk{}structureOK} --- from the fact that citations point
  strictly backwards, which needs only sortedness;
\item and rule correctness, \isa{hlEmitDerivationFuel\ubrk{}ruleOK}, over all
  twenty-one rules.
\end{itemize}

Three things carry the last of these. First,
\isa{hlFitchToLemmon\ubrk{}references}: the partial environment the erasure held when
it computed a line gives, at every line that rule may cite, the same answer as
reading the whole finished proof back --- later lines cannot interfere because
their numbers are larger and citations point backwards. Second,
\isa{hlEmitDerivationFuel\ubrk{}deps}: an emitted line's dependencies are exactly the
environment image of its subderivation's open assumptions. The load-bearing
step there is \isa{hlEnvImage\ubrk{}drop}, which says that because a box head is
freshly allocated, subtracting it removes exactly the discharged assumption and
never an older one --- the same hazard class as the \isa{boxHeadError} defect of
\S\ref{sec:direct}.

Third, the eigenconstant obstruction, which was real. At the level of the
derivation, \isa{hlForallIntroStep} and \isa{hlExistsElimStep} state freshness
against the \emph{derivation's} open assumptions; \isa{hlRuleOK} states it against the
\emph{proof's} assumption lines. That gap is closed, given the dependency
equation, by \isa{hlAssumptionConstants\ubrk{}envImage} together with
\isa{hlReferencedConstants\ubrk{}envImage}.

\subsection{The scope check}\label{sec:scopecheck}

\isa{hlFitchCorrect} requires more than the above: the eigenconstant conditions
must hold against the assumptions genuinely in Fitch \emph{scope}, which is
\isa{hlCorrectG (hlScopeSrc F)}. This does not follow from what precedes it. A
line's scope contains every root premise and every enclosing box head, so it is
in general \emph{larger} than that line's dependency set, and the antitonicity
lemmas in \thy{LF\ubrk{}HLW\ubrk{}Rule\ubrk{}Transfer.thy} run the wrong way.

\isa{hlEmitDerivationFuel\ubrk{}scopeOK} proves it by a second induction over the
emitter carrying three further invariants: that the assumption constants of the
scope set are contained in the constants of the emitter's own \isa{scope} list;
that every number in the scope set is below the next free line, so a freshly
allocated box head is never already in scope; and \isa{hlScopeBelow}, that every
constant the scope holds is shorter than the name the next repair will mint.
Nineteen of the twenty-one rules reduce immediately, since only
\isa{HL\ubrk{}ForallIntro} and \isa{HL\ubrk{}ExistsElim} consult the scope at all. What
makes the remaining two go through is discussed in \S\ref{sec:found}.

The freshness bookkeeping is in \thy{LF\ubrk{}HLW\ubrk{}Fresh.thy}:
\isa{hlNamesBelow base cnt d} says every constant in $d$ is shorter than the
name the next repair will mint, which survives each repair because a repair adds
one name of exactly that length and increments the counter. That is what the
counter in \isa{hlFreshIndex} is for, and the invariant transfers to every
subderivation by \isa{hlEmitDerivationFuel\ubrk{}count\ubrk{}le} and
\isa{hlDerivationConstants\ubrk{}sub}.

\subsection{The construction theorem}

\begin{theorem*}[\isa{hlPaperCorrect\ubrk{}toFitch}, \thy{LF\ubrk{}HLW\ubrk{}Construct.thy}]
For every nonempty proof $P$ satisfying \isa{hlPaperCorrect} there are a
route and an $F$ with
\begin{gather*}
  \isa{hlLemmonToFitchChecked}\;P = \mathrm{Inr}(\mathit{route},F),
  \qquad \isa{hlFitchCorrect}\;F, \\
  \isa{hlFitchConclusion}\;F = \isa{hlConclusion}\;P, \qquad
  \isa{hlConclusion}(\delta_H F) = \isa{hlConclusion}\;P, \\
  \mathrm{set}(\isa{hlFitchPremises}\;F) \subseteq
    \mathrm{set}(\isa{hlOpenPremises}\;P).
\end{gather*}
\end{theorem*}

\noindent
The route is existential because the checked operation tries the direct
algorithm first and falls back to the tree construction; since Conjecture~27 is
false, no single direct algorithm can be claimed, and the theorem is about the
translation as a whole. This is strictly stronger than the earlier
\isa{hlViaTree\ubrk{}accepted} and \isa{hlLemmonToFitchChecked\ubrk{}accepted}, which
guarantee any \emph{particular} accepted output but cannot rule out the
translator rejecting a good source.

Assembling it required discharging the inductions' hypotheses at the root:
\isa{hlEmitCtx} bundles the requirements that the layout environment cover every
open assumption, that each environment line be its own sole dependency, and that
each carry that assumption's formula. All three come from the premise layout.
Premise closure (\isa{hlFitchPremiseClosed}) follows because the conclusion
line's dependencies are the environment image of the tree's open assumptions,
every one of which is a premise line; the case of a derivation that is a single
premise leaf is separate, since then the last line is itself a premise line.

\subsection{One hypothesis}

The construction theorem needs one side condition on the derivation tree: no
assumption label may carry two different formulas
(\isa{hlAssumptionsFunctional}). \isa{hlEnvironmentLine} returns the
\emph{first} environment entry with a given label, so without this the
environment need not describe the tree's open assumptions at all. It is
automatic for anything arising from a proof --- a label is a line number and a
line carries one formula --- and \isa{hlToDerivation\ubrk{}functional} discharges it,
since the unfolding places the open assumptions inside
\isa{hlSourceAssumptions}, which is functional because \isa{hlLookupLine} is. It
therefore does not appear in the theorem above, but it is a genuine condition on
arbitrary \isa{hl\ubrk{}derivation} values.

\section{The impossibility results}\label{sec:impossibility}

\subsection{Three obstructions}

The paper's Theorem~10 says some correct Lemmon proofs have no positional Fitch
image, and Conjectures~27 and~28 propose weaker forms that might survive. All
three are refuted, and there are exactly three obstructions, which the
development keeps distinct because they are genuinely different phenomena.

\begin{center}
\begin{tabular}{@{}lll@{}}
\toprule
Obstruction & Witness & Refutes \\
\midrule
Discharges do not nest    & \isa{ex11} & Theorem 10 \\
A line is trapped in a box & \isa{ex12} & Theorem 10 \\
Two discharges share a last line & \isa{c27} & Conjectures 27, 28 \\
\bottomrule
\end{tabular}
\end{center}

\noindent
The third is worth emphasising, because it is easy to mistake for one of the
others. In \isa{c27} the two boxes nest perfectly and no line is trapped; what
fails is that both conditional proofs discharge \emph{from the same line}, one
releasing assumption 1 and the other assumption 2, so a Fitch proof would need
two boxes ending at the same place.

\subsection{At the compact types}

\isa{theorem\ubrk{}10\ubrk{}ex11}, \isa{theorem\ubrk{}10\ubrk{}ex12} and \isa{theorem\ubrk{}10} in
\thy{LF\ubrk{}Theorem10.thy} are genuine quantified nonexistence statements, not
observations that one algorithm returns an error. They rest on the positional
apparatus of \thy{LF\ubrk{}Span.thy} and \thy{LF\ubrk{}Positional.thy}: the compact
flatten records, for each line, the \emph{path} of boxes containing it, and the
arguments are prefix reasoning over those paths, with \isa{subs\ubrk{}span} as the
central lemma. \isa{conjecture\ubrk{}27\ubrk{}false} and \isa{conjecture\ubrk{}28\ubrk{}false} in
\thy{LF\ubrk{}Conjecture27.thy} do the same for the conjectures, the first
quantifying over all injective renumberings, so it covers the permitted
permutations rather than one of them.

\subsection{At the exact types}\label{sec:bridge}

Carrying these to the exact types looked like the hardest part of the project,
and the two obvious routes both have problems.

A \emph{bridge} between the layers would transport the results, but it must run
from exact to compact --- to show no exact image exists, assume one and map it
down --- and in that direction it cannot be total: \isa{HL\ubrk{}MT},
\isa{HL\ubrk{}LEM}, \isa{HL\ubrk{}PropTaut} and \isa{HL\ubrk{}QN} have no compact counterpart,
and the last is an arbitrary propositional-consequence rule. (The other
direction is nearly total: \isa{BotI} maps to \isa{HL\ubrk{}PropTaut [i,j]},
\isa{Reit} to \isa{HL\ubrk{}PropTaut [i]}, and the split eliminations collapse onto
the single exact ones.) A witness-local bridge would have sufficed, since
\isa{lineMatches} pins an image's rules to the source's, but that still leaves
relating two independently written well-formedness predicates. \emph{Porting}
the positional apparatus was the alternative, at the cost of most of a
2{,}000-line theory.

Neither was needed. The exact layer already carries the same content in a
different form, because \isa{hlFitchNestingFrom} admits a discharge citation
only when the pair is a box that has closed at the citing line's own level, and
tracks a visible set that a closed box's interior never re-enters. Three
structural facts follow directly, by the induction that
\isa{hlBoxSpans} itself follows:

\begin{description}
\item[\isa{hlBoxSpans\ubrk{}nested}] Box spans never cross: if $a < a' \le c$ for
  spans $(a,c)$ and $(a',c')$, then $c' \le c$.
\item[\isa{hlFitchNestingFrom\ubrk{}no\ubrk{}cite\ubrk{}into\ubrk{}box}] A line numbered past a
  box's last line cannot cite a number inside it. The induction carries the
  invariant that the visible set holds only numbers earlier than the fragment.
\item[\isa{hlBoxSpans\ubrk{}last\ubrk{}lt}, \isa{hlBoxSpans\ubrk{}same\ubrk{}last}] A box ends
  strictly before the proof containing it, so two distinct boxes never share a
  last line.
\end{description}

\noindent
Each obstruction is then one of these. \isa{hl\ubrk{}theorem\ubrk{}10\ubrk{}ex11} closes at once:
the two conditional proofs would need boxes spanning 1 to 3 and 2 to 4, and
$1 < 2 \le 3$ forces $4 \le 3$. \isa{hl\ubrk{}theorem\ubrk{}10\ubrk{}ex12} is a little longer.
\isa{hlSharedDischarge\ubrk{}no\ubrk{}image} is again immediate, and
\isa{hl\ubrk{}conjecture\ubrk{}27\ubrk{}false} follows from it by
\isa{hlSharedDischarge\ubrk{}permute}: sharing a discharge is a property of the
proof rather than of the page, so it survives every injective renumbering.

\section{What the formalisation found}\label{sec:found}

\subsection{Corrections to the paper}

\begin{enumerate}
\item \textbf{Theorem 4's converse is false.} Retain totality and the forward
  implication; the converse is refuted by the paper's own Theorem~22 witness.
\item \textbf{Proposition 5 needs a corrected notion of scope.} Literal
  Definition~2 omits outer premises from the assumptions in scope, which makes
  Proposition~5 false already at a premise line. \isa{scopeOf} includes them.
\item \textbf{Definition 17's graph orientation must be stated explicitly.} The
  paper writes cited$\to$citing edges, so unfolding from the conclusion
  traverses the reverse relation. The code's backward recursion is the intended
  algorithm, not a literal traversal of the written edges.
\item \textbf{Conjecture 26 needs both eigenconstant repairs.} The original
  repairs universal introduction only. Existential elimination's witness
  carries the same kind of freshness condition and needs the same treatment;
  the proved compact result is a corrected construction theorem rather than the
  original statement.
\end{enumerate}

\noindent
A fifth item is a matter of interpretation rather than error: the paper's later
remarks about normalisation do not amount to a retract on raw proof objects,
since unfolding can duplicate lines, insert reiterations, rename constants and
omit lines irrelevant to the conclusion. The revised manuscript withdraws that
claim.

\subsection{Defects in the supplied Haskell}

Fourteen findings in the supplied source are recorded in
\thy{HASKELL\ubrk{}AUDIT.md}, together with two in the code Isabelle generates.
Most of the fourteen are partiality:
omitted \isa{Boolean} cases in six \thy{ProofTypes} functions, \isa{Iff} missing
from \isa{ModelSemantics.eval} and from \thy{PropDNF}, six justification
constructors unhandled in \thy{JsonToPipe}. Two are more serious.

\paragraph{A structural check that does not check order.}
\isa{LemmonChecker.checkStructure} compares cited line \emph{numbers} with the
current number rather than checking physical proof order. A first line numbered
10 can cite a later line numbered 5 and evade the intended acyclicity check.
This is why \isa{hlVerifiedCorrect} exists as a separate predicate above the
faithful \isa{hlCorrect}.

\paragraph{A soundness bug.}
\isa{FitchTypes.fitchWellFormed} checks scope and citations but not rule
correctness, eigenvariables, the existence of the subproofs a discharge refers
to, or a root-level final conclusion. The last of these is the soundness bug. A proof consisting of one bare subproof therefore has no
premises and reports the subproof's \emph{assumption} as its conclusion --- so
it ``proves'' an arbitrary formula from nothing. The witness is
\isa{hlOpenAssumptionFitch}, and \isa{hlOpenAssumptionFitch\ubrk{}fails\ubrk{}only\ubrk{}root\ubrk{}scope}
shows by evaluation that it passes every other condition, including the
scope-based rule check. The faithful \isa{hlFitchWellFormed} reproduces the
defect deliberately; the condition was added to \isa{hlConcludesAtTop}, which
the layers above require. The same repair in the compact layer turned a refuted
soundness claim into the proved \isa{fitch\ubrk{}semantic\ubrk{}soundness}.

\paragraph{A translation bug.}
\isa{FitchConvert.toDerivation} computes one global set of discharged
assumptions. If an assumption is discharged on a line that is dead with respect
to the selected conclusion, the tree still labels that assumption as discharged
rather than as a premise; the shared-discharge proof then loses an open premise
and the emitter cites line~0. Isabelle corrects this with
\isa{hlClassifyAssumptions}, whose bound set follows actual derivation
ancestors.

\subsection{Two things the main theorem depends on}

Two design decisions in the original system turned out to be load-bearing in
ways that are not obvious from reading it.

\paragraph{The witness inference draws candidates from the premise.}
For the scope check of \S\ref{sec:scopecheck} to go through, an inferred
eigenconstant must lie in the set the repair empties of scope constants ---
namely the constants of the conclusion that are missing from the goal. It does,
and the reason is in the checker: \isa{hlWitnessLists} takes its candidate
constants from \isa{hlConstantsInFormula q}, the constants of the
\emph{premise}. So an inferred witness always occurs in the premise and, having
been abstracted away, never in the conclusion. Had the candidates been drawn
from anywhere else, the scope check would not follow, and this would be a gap
rather than a theorem. The formal statements are
\isa{hlForallIntroStep\ubrk{}witness} and \isa{hlExistsElimStep\ubrk{}witness}.

\paragraph{\isa{hlConcludesAtTop} is not presentational.}
The requirement that a subproof's body end with a line rather than another
subproof looks like a matter of style. Without it, a box whose body ends with a
box \emph{inherits that box's last line}, and \isa{hlBoxSpans\ubrk{}same\ubrk{}last} is
false --- so Conjectures~27 and~28 would not be refutable by this route. The
condition the nesting checker imposes on every body is exactly what excludes
that case. It is also, as above, what repairs the soundness bug in
\isa{fitchWellFormed}. One condition is doing two quite different jobs.

\section{Validation}\label{sec:validation}

\paragraph{Proof checking.} The complete development processes from Isabelle's
plain \isa{HOL} image in about six minutes:
\isa{isabelle build -d . -e Lemmon\ubrk{}Fitch}. A source scan finds no \isa{sorry},
no \isa{oops}, no axiomatizations, no oracles and no \isa{quick\ubrk{}and\ubrk{}dirty}
settings, and every maintained theory is explicitly imported by the cap, so
processing the cap checks all of it. Of the 1{,}234 results, 124 are settled by
evaluation on concrete witnesses --- the counterexamples and the regression
suites --- and the rest by proof. Isabelle 2025-2; only \isa{HOL} and
\isa{HOL-Library} are required, and no Archive of Formal Proofs entry.

\paragraph{Differential testing.} The development generates a Haskell kernel
and an SML kernel by code extraction. The generated checker is compared against
the original handwritten one, per line and through its public API rather than
through internals: 9{,}874 comparisons --- the supplied 14-proof corpus,
twelve hand-written proofs covering the eight rules that corpus never
exercises, and 1{,}728 systematic mutants --- with 0 verdict disagreements, of
which 2{,}706 exercise the rejection path; 860 of 860 rejection messages
reproducing the original text byte for byte; all 80 cases of the supplied
regression suite passing with the generated kernel substituted; 26 of 26 public
API assertions. The messages come from \isa{hlLineError}, which Isabelle proves
reports no error exactly when \isa{hlLineOK} holds, so the diagnostics cannot
alter which proofs are accepted. This is testing, not proof --- the two can never
be proved equivalent, since the Haskell has no formal semantics here --- but it
is the check that distinguishes a faithful reconstruction from a careful
lookalike.

\paragraph{Regression.} \thy{LF\ubrk{}HLW\ubrk{}Regression.thy} and
\thy{LF\ubrk{}HLW\ubrk{}Delta.thy} evaluate the checker on the configurations that
previously escaped it: closed-box escapes, nonexistent discharge pairs,
cross-level subproof citations, malformed placements, and the tree route on
examples of all twenty-one rules.

\section{What is not established}\label{sec:limits}

\begin{itemize}
\item \textbf{Whether \isa{hlCorrect} alone is semantically sound.} Soundness is
  proved of the stronger \isa{hlVerifiedCorrect}. Settling this needs a witness
  rather than an argument from the definitions; the \isa{fitchWellFormed} bug
  above is the standing reason for caution.
\item \textbf{Whether the seven direct-translation checks are necessary.}
  Sufficiency is proved. All six witnesses considered are genuinely imageless,
  but six witnesses are not all sources.
\item \textbf{Conjecture 28 under unrestricted citation rewriting.} The
  refutation fixes an explicit formula-, rule- and citation-preserving notion of
  image, and allows auxiliary lines
  (\isa{conjecture\ubrk{}28\ubrk{}with\ubrk{}auxiliaries\ubrk{}false}). A broader notion in which
  source citations may be redirected to auxiliary derivations is untouched, and
  the original gives no formal definition to appeal to.
\item \textbf{Graph-theoretic presentation.} Definitions 14--15 and
  Propositions 18 and 20 are formalised in their operational content ---
  acyclicity as a descending citation measure, scope as nested datatypes --- not
  as a separate general graph theory with an isomorphism theorem.
\item \textbf{The application layer.} The web application has not been adapted
  to the generated Aeson and parser interface. This is interface engineering
  rather than formalisation.
\end{itemize}

\appendix

\section{Theory map}

\begin{longtable}{@{}lrr l@{}}
\toprule
Theory & Lines & Results & Subject \\
\midrule
\endhead
\thy{LF\ubrk{}Formula} & 319 & 23 & compact formulas, renaming, abstraction \\
\thy{LF\ubrk{}Lemmon} & 449 & 12 & compact Lemmon proofs and checker \\
\thy{LF\ubrk{}Fitch} & 346 & 8 & compact Fitch syntax, scope paths \\
\thy{LF\ubrk{}Delta} & 463 & 38 & erasure; Theorem 4, Propositions 5 and 7 \\
\thy{LF\ubrk{}Direct} & 314 & 5 & the direct algorithm and its seven checks \\
\thy{LF\ubrk{}Derivation} & 583 & 37 & derivation trees; Lemma 23, Corollary 24 \\
\thy{LF\ubrk{}Unfold} & 561 & 19 & unfolding; L1, Proposition 18 \\
\thy{LF\ubrk{}Render} & 236 & 1 & printing \\
\thy{LF\ubrk{}Examples} & 403 & 24 & the paper's examples; Theorem 22 \\
\thy{LF\ubrk{}WellFormed} & 1906 & 91 & L2, L3 \\
\thy{LF\ubrk{}Faithful} & 712 & 20 & source-to-tree faithfulness \\
\thy{LF\ubrk{}Check} & 2465 & 80 & L4, rule transfer \\
\thy{LF\ubrk{}Conjecture} & 228 & 9 & compact construction theorem \\
\thy{LF\ubrk{}Span} & 577 & 28 & box spans \\
\thy{LF\ubrk{}Positional} & 2014 & 102 & positional analysis of the direct route \\
\thy{LF\ubrk{}Semantics} & 854 & 44 & compact semantics and soundness \\
\thy{LF\ubrk{}Theorem10} & 246 & 10 & Theorem 10, compact \\
\thy{LF\ubrk{}Conjecture27} & 719 & 46 & Conjectures 27 and 28, compact \\
\midrule
\thy{LF\ubrk{}HLW\ubrk{}Formula} & 339 & 2 & exact formulas \\
\thy{LF\ubrk{}HLW} & 545 & 3 & exact proofs; the twenty-one rule checks \\
\thy{LF\ubrk{}HLW\ubrk{}Canonical} & 158 & 12 & \isa{hlPaperCorrect}; Proposition 7 \\
\thy{LF\ubrk{}HLW\ubrk{}Fitch} & 334 & 5 & exact Fitch syntax; nesting \\
\thy{LF\ubrk{}HLW\ubrk{}Scope} & 291 & 7 & scope records; \isa{hlFitchCorrect} \\
\thy{LF\ubrk{}HLW\ubrk{}Semantics} & 2814 & 92 & exact semantics; \isa{hlRuleOK\ubrk{}sound} \\
\thy{LF\ubrk{}HLW\ubrk{}Fitch\ubrk{}Semantics} & 41 & 2 & Fitch soundness \\
\thy{LF\ubrk{}HLW\ubrk{}DNF} & 137 & 3 & propositional normal forms \\
\thy{LF\ubrk{}HLW\ubrk{}Translate} & 299 & 4 & exact direct translation \\
\thy{LF\ubrk{}HLW\ubrk{}Unfold} & 663 & 3 & exact unfolding and emitter \\
\thy{LF\ubrk{}HLW\ubrk{}Translation\ubrk{}Proofs} & 354 & 18 & L1; conditional guarantees \\
\thy{LF\ubrk{}HLW\ubrk{}Derivation} & 251 & 10 & exact derivation trees \\
\thy{LF\ubrk{}HLW\ubrk{}Faithful} & 559 & 15 & G1: source to correct tree \\
\thy{LF\ubrk{}HLW\ubrk{}Renaming} & 390 & 47 & renaming foundations \\
\thy{LF\ubrk{}HLW\ubrk{}Witnesses} & 199 & 13 & witness inference \\
\thy{LF\ubrk{}HLW\ubrk{}Quantifier\ubrk{}Renaming} & 149 & 12 & renaming and quantifiers \\
\thy{LF\ubrk{}HLW\ubrk{}Derivation\ubrk{}Renaming} & 172 & 18 & G2: renaming preserves rules \\
\thy{LF\ubrk{}HLW\ubrk{}Rule\ubrk{}Transfer} & 320 & 12 & local transfer; antitonicity \\
\thy{LF\ubrk{}HLW\ubrk{}Layout} & 222 & 27 & G3: allocation \\
\thy{LF\ubrk{}HLW\ubrk{}Fuel} & 66 & 3 & fuel adequacy \\
\thy{LF\ubrk{}HLW\ubrk{}Conclusion} & 75 & 8 & conclusion preservation \\
\thy{LF\ubrk{}HLW\ubrk{}Premise\ubrk{}Layout} & 161 & 20 & premise placement \\
\thy{LF\ubrk{}HLW\ubrk{}Nesting} & 944 & 41 & T1: structural well-formedness \\
\thy{LF\ubrk{}HLW\ubrk{}Fresh} & 501 & 30 & the repair and its name bound \\
\thy{LF\ubrk{}HLW\ubrk{}Erasure} & 3793 & 95 & T2: rule transfer through $\delta_H$ \\
\thy{LF\ubrk{}HLW\ubrk{}Construct} & 3408 & 55 & T3: the construction theorem \\
\thy{LF\ubrk{}HLW\ubrk{}Theorem10} & 996 & 35 & impossibility at exact types \\
\thy{LF\ubrk{}HLW\ubrk{}Examples} & 317 & 21 & the paper's examples, exact \\
\thy{LF\ubrk{}HLW\ubrk{}Delta} & 106 & 13 & erasure regressions \\
\thy{LF\ubrk{}HLW\ubrk{}Regression} & 69 & 7 & all-rule regressions \\
\thy{LF\ubrk{}HLW\ubrk{}Report} & 304 & 4 & diagnostic reporting \\
\bottomrule
\end{longtable}

\section{Claim to theorem}

\begingroup\small
\begin{longtable}{@{}>{\raggedright\arraybackslash}p{0.26\textwidth}
                   >{\raggedright\arraybackslash}p{0.31\textwidth}
                   >{\raggedright\arraybackslash}p{0.33\textwidth}@{}}
\toprule
Paper item & Compact & Exact \\
\midrule
\endhead
Theorem 4, forward & \isa{theorem\ubrk{}4\ubrk{}forward} & --- \\
Theorem 4, converse & \isa{theorem\ubrk{}4\ubrk{}backward\ubrk{}fails} & \isa{hl\ubrk{}theorem\ubrk{}4\ubrk{}backward\ubrk{}fails} \\
Proposition 5 & \isa{proposition\ubrk{}5} & --- \\
Corollary 6 & \isa{corollary\ubrk{}6} & --- \\
Proposition 7 & \isa{proposition\ubrk{}7} & \isa{hl\ubrk{}proposition\ubrk{}7} \\
Theorem 10 & \isa{theorem\ubrk{}10} & \isa{hl\ubrk{}theorem\ubrk{}10} \\
Proposition 18 & \isa{proposition\ubrk{}18} & --- \\
Theorem 22 & \isa{theorem\ubrk{}22} & \isa{hl\ubrk{}thm22\ubrk{}*} \\
Lemma 23 & \isa{renaming} & \isa{hlDerivationOK\ubrk{}rename} \\
Corollary 24 & \isa{renaming\ubrk{}open} & \isa{hlRenameDerivation\ubrk{}open\ubrk{}unchanged} \\
Conjecture 26 (corrected) & \isa{conjecture\ubrk{}26} & \isa{hlPaperCorrect\ubrk{}toFitch} \\
Conjecture 27 & \isa{conjecture\ubrk{}27\ubrk{}false} & \isa{hl\ubrk{}conjecture\ubrk{}27\ubrk{}false} \\
Conjecture 28 & \isa{conjecture\ubrk{}28\ubrk{}false} & \isa{hl\ubrk{}conjecture\ubrk{}28\ubrk{}false} \\
L1 & \isa{L1} & \isa{hlToDerivation\ubrk{}total} \\
L2, L3 & \isa{L2\ubrk{}L3} & \isa{hlClassified\brk{}Derivation\brk{}ToFitch\ubrk{}nested\brk{}WellFormed} \\
L4 & \isa{L4\ubrk{}derivation} & \isa{hlEmitDerivationFuel\ubrk{}ruleOK}, \isa{\dots\ubrk{}scopeOK} \\
Soundness & \isa{lemmon\ubrk{}semantic\ubrk{}soundness} & \isa{hlRuleOK\ubrk{}sound} \\
\bottomrule
\end{longtable}
\endgroup

\end{document}
